\section*{Chapter \ref{ch:g11:infrared} --- Infrared Spectroscopy}

\begin{solution}{exo:g11:infrared:1}
$\lambda = 1/3300 = \qty{3.03e-4}{cm} = \qty{3.03}{\micro m}$;
$\lambda = 1/1050 = \qty{9.52e-4}{cm} = \qty{9.52}{\micro m}$. The band at
\qty{3300}{cm^{-1}}, of larger \omterm{def:g11:infrared:wavenumber}{wavenumber}, belongs to the faster
vibration.
\end{solution}

\begin{solution}{exo:g11:infrared:2}
$\qty{10}{\micro m} = \qty{1.0e-3}{cm}$, so $\sigma = 1/\num{1.0e-3} =
\qty{1000}{cm^{-1}}$.
\end{solution}

\begin{solution}{exo:g11:infrared:3}
The \ce{C=O} of a \omterm{def:g11:functional-groups:carbonyl}{ketone}; a hydrogen-bonded \omterm{def:g11:functional-groups:alcohol}{alcohol} \ce{O-H}; \ce{C-H}
bonds of a carbon chain.
\end{solution}

\begin{solution}{exo:g11:infrared:4}
The band at \qty{1715}{cm^{-1}}: the \ce{C=O} bond.
\end{solution}

\begin{solution}{exo:g11:infrared:5}
The instrument measures the light that passes through the sample; where
a bond absorbs, less light passes and the curve dips. $T =
\qty{100}{\percent}$ means that none of the light of that \omterm{def:g11:infrared:wavenumber}{wavenumber} is
absorbed.
\end{solution}

\begin{solution}{exo:g11:infrared:6}
The first has a \ce{C=O} band at \qty{1715}{cm^{-1}} and no \ce{O-H}
band. It is a \omterm{def:g11:functional-groups:carbonyl}{ketone}: \ce{CH3-CO-CH2-CH3}, butanone (also called
butan-2-one). The second has an \ce{O-H} band and no \ce{C=O}: an
\omterm{def:g11:functional-groups:alcohol}{alcohol}, for example butan-1-ol.
\end{solution}

\begin{solution}{exo:g11:infrared:7}
The \ce{O-H} groups of the acid are very strongly hydrogen-bonded (the
\omterm{def:g7:atoms-and-molecules:molecule}{molecules} pair up, each \ce{O-H} bonded to the \ce{C=O} of the other),
and every \omterm{def:g7:atoms-and-molecules:molecule}{molecule} is bonded a little differently: the band is loosened
further and spread over a very wide range.
\end{solution}

\begin{solution}{exo:g11:infrared:8}
In the liquid (A) the \ce{O-H} band is broad, near
\qty{3350}{cm^{-1}}; in the gas (B), where the \omterm{def:g7:atoms-and-molecules:molecule}{molecules} are far apart
and form no \omterm{def:g11:polarity-and-cohesion:hydrogen-bond}{hydrogen bonds}, it is a sharp band near
\qty{3666}{cm^{-1}}.
\end{solution}

\begin{solution}{exo:g11:infrared:9}
Hardly: near \qty{1730}{cm^{-1}} for the \omterm{def:g11:functional-groups:carbonyl}{aldehyde} and
\qty{1715}{cm^{-1}} for the \omterm{def:g11:functional-groups:carbonyl}{ketone}, close values. Comparison of the
whole spectrum, \omterm{def:g11:infrared:band}{fingerprint region} included, with reference spectra
would decide, or another technique.
\end{solution}

\begin{solution}{exo:g11:infrared:10}
An \omterm{def:g11:functional-groups:carboxylic-acid}{ester} (\ce{C=O} near \qty{1735}{cm^{-1}}, strong \ce{C-O}, no
\ce{O-H}), for example ethyl ethanoate. Ethanoic acid is \ce{C2H4O2}, not
\ce{C4H8O2}. Butanoic acid has the right formula but would show the very
broad \ce{O-H} band of acids: excluded.
\end{solution}

\begin{solution}{exo:g11:infrared:11}
Propan-1-ol is an \omterm{def:g11:functional-groups:alcohol}{alcohol}: like A. Propanoic acid is a \omterm{def:g11:functional-groups:carboxylic-acid}{carboxylic acid}:
like D.
\end{solution}

\begin{solution}{exo:g11:infrared:12}
The broad \ce{O-H} band near \qty{3350}{cm^{-1}} shrinks while a strong
\ce{C=O} band near \qty{1715}{cm^{-1}} grows. The reaction is complete
when the \ce{O-H} band has disappeared and the \ce{C=O} band no longer
grows.
\end{solution}

\begin{solution}{exo:g11:infrared:13}
Butanoic acid shows a very broad \ce{O-H} band from 2500 to
\qty{3300}{cm^{-1}}, the \omterm{def:g11:functional-groups:carboxylic-acid}{ester} none; the acid's \ce{C=O} lies near
\qty{1710}{cm^{-1}}, broader, the \omterm{def:g11:functional-groups:carboxylic-acid}{ester}'s near \qty{1735}{cm^{-1}},
sharp, with a strong \ce{C-O} band near \qty{1240}{cm^{-1}}. The acid has
an \ce{O-H} group, hydrogen-bonded, which also slightly loosens its
\ce{C=O}.
\end{solution}

\begin{solution}{exo:g11:infrared:14}
Leftover ethanol (the \omterm{def:g11:functional-groups:alcohol}{alcohol} the \omterm{def:g11:functional-groups:carboxylic-acid}{ester} is made from) or water (from
the \omterm{def:g4:air-a-mixture-of-gases:air}{air} or the washing): both have hydrogen-bonded \ce{O-H} groups. The
sample can be compared with a reference spectrum, dried, distilled, or
its boiling temperature measured.
\end{solution}

\begin{solution}{exo:g11:infrared:15}
A \omterm{def:g10:lewis-and-shape:multiple-bond}{double bond} is a stiffer spring than a single bond between the same
\omterm{def:g7:atoms-and-molecules:atom}{atoms}: it vibrates faster, at a higher \omterm{def:g11:infrared:wavenumber}{wavenumber} (1715 against
\qty{1050}{cm^{-1}}). The \ce{O-H}, \ce{N-H} and \ce{C-H} bonds all
involve a hydrogen \omterm{def:g7:atoms-and-molecules:atom}{atom}, the lightest of all: light balls on a spring
vibrate fast, hence \omterm{def:g11:infrared:wavenumber}{wavenumbers} near 2900--\qty{3600}{cm^{-1}}.
\end{solution}

\begin{solution}{pb:g11:infrared:1}
\textbf{1.} \ce{CH3-CH2-OH}; \ce{CH3-CO-CH3}; \ce{CH3-COOH}.

\textbf{2.} \ce{C2H6O}; \ce{C3H6O}; \ce{C2H4O2}.

\textbf{3.} Hydroxyl; carbonyl (\omterm{def:g11:functional-groups:carbonyl}{ketone}); carboxyl.

\textbf{4.} Ethanol: \ce{O-H} and \ce{C-H}. Propanone: \ce{C-H} and
\ce{C=O}. Ethanoic acid: \ce{O-H}, \ce{C-H} and \ce{C=O}.

\textbf{5.} No \ce{C=O} band, a broad \ce{O-H} band: ethanol.

\textbf{6.} A very broad \ce{O-H} band from 2500 to \qty{3300}{cm^{-1}}
and a strong \ce{C=O} near \qty{1710}{cm^{-1}}: ethanoic acid.

\textbf{7.} Propanone: the very strong \ce{C=O} band at
\qty{1715}{cm^{-1}}, with no \ce{O-H}.

\textbf{8.} Probably (vinegar, \omterm{def:g6:solutions-and-solubility:solute}{solvent}, \omterm{def:g11:functional-groups:alcohol}{alcohol}), but sniffing unknown
liquids is dangerous; the spectrum is safe and certain.

\textbf{9.} The acid's \ce{O-H} groups are more strongly hydrogen-bonded
(paired \omterm{def:g7:atoms-and-molecules:molecule}{molecules}), so the band is lower and much wider.

\textbf{10.} \ce{CH3-O-CH3}, \ce{C2H6O}; it has no \ce{O-H}, so no
\ce{O-H} band.

\textbf{11.} The \ce{C=O} band: near \qty{1730}{cm^{-1}} for propanal,
near \qty{1715}{cm^{-1}} for propanone.

\textbf{12.} The \ce{O-H} band: an \omterm{def:g11:functional-groups:carboxylic-acid}{ester} has no \ce{O-H}. Its \ce{C=O}
would lie near \qty{1735}{cm^{-1}}.

\textbf{13.} \omterm{def:g11:organic-skeletons:isomer}{Isomers} share a formula; their \omterm{def:g11:functional-groups:functional-group}{functional groups} differ,
and the spectrum shows the groups.

\textbf{14.} \qty{1715}{cm^{-1}}.

\textbf{15.} $1715 \times 100 = \qty{1.715e5}{m^{-1}}$.

\textbf{16.} $\lambda = 1/\num{1.715e5} = \qty{5.83e-6}{m} =
\qty{5.83}{\micro m}$.

\textbf{17.} No: visible light runs from about 0.4 to
\qty{0.75}{\micro m}. It is infrared radiation.

\textbf{18.} $\num{1.05e5}\ \si{m^{-1}}$, so $\lambda = \qty{9.52e-6}{m}
= \qty{9.52}{\micro m}$.

\textbf{19.} The \ce{C=O} bond (higher \omterm{def:g11:infrared:wavenumber}{wavenumber}, shorter wavelength).

\textbf{20.} About \qty{5.8}{\micro m}.
\end{solution}
